\section{Sum-of-Squares Certificate of Expansion in the Johnson Graph}\label{sec:fourier}
In this section we give an \sos proof of Theorem~\ref{thm:structure-johnson}.
Our proof follows the same lines as~\cite{KMMS}, except that we implement each step carefully by a low-degree SoS proof. For convenience, we restate 
Theorem~\ref{thm:structure-johnson} below.
\begin{theorem}[Expansion Theorem \ref{thm:structure-johnson} for Johnson Graphs  restated]\label{thm:structure-johnson-restated}
For all $\alpha \in (0,1)$, all integers $\ell \geq 1/\alpha$ and $n \geq \ell$, the following holds: Let $J_{n,\ell,\alpha}$ be the $\alpha$-noisy Johnson graph. For every constant $\gamma \in (0,1)$ and positive integer $r \leq O(\ell)$, every function $F: V(J) \rightarrow \R$ that is $(r,\gamma)$-pseudorandom has high expansion:
\begin{enumerate}
\item $\{F(X) \in [0,1]\}_{X \in V(J)}  \vdash_{O(1)}\, \\
\ip{F,LF} \geq \delta(1 - (1-\alpha)^{r+1})(1- \gamma^{1/3}\exp(r)) - \sum_{j = 0}^r \frac{ c_j\ell^{j}}{\gamma}\E_{a \sim {[n] \choose j}}[q_a(F)({\delta(F^{ 2}|_a) - \gamma)}]+ B(F)$
\item $\{F(X) \in [0,1]\} \,\,\, \vdash_{2} \,\,\, 0 \leq q_a(F) \leq 1,$
\end{enumerate}
where for all $j \leq r$, $c_j$'s are positive constants of size at most $\exp(r)$, for all 
$a \subseteq [n]$ of size $j$, $q_a(F)$ is a degree $2$  polynomials, and $B(F)=\frac{4}{3}\E_X[(F^{ 3} - F)\Pi_{\geq \lambda_{r}}F]$ for $\lambda_r = (1-\alpha)^r$. Here, $\Pi_{\geq \lambda}$ is the projection operator on the space of spanned by eigenvectors of$J$ of eigenvalues at least $\lambda$.
\end{theorem}

It will be more convenient for us to move to a closely related Cayley graph, which we denote by $C_{n,\ell,\alpha}$. This graph is essentially the same as the Johnson graph, albeit viewed as a product domain. In Theorem~\ref{thm:structure-johnson-cayley}, we state an analogous structure theorem for ``permutation-invariant'' sets on $C_{n,\ell,\alpha}$ that we show in Theorem~\ref{thm:structure-johnson-cayley}, and it is easy to derive Theorem~\ref{thm:structure-johnson-restated} above from it.



\paragraph{Notation.} 
We use $[n]$ to denote the set $\{0,\ldots,n-1\}$, and also the group $(\Z/n\Z)$, the natural numbers modulo $n$. 
Generally, when we take a set $S$ and raise it to a positive integer power $\ell$, we mean the set of all ordered multisets of elements of $S$ of size $\ell$. We use $\Chi_t$ for $t \in [n]$ to denote the characters of the group $\Z/n\Z$ (or the eigenvectors of the $n$-cycle), where $\Chi_t: [n] \rightarrow \mathbb{C}$ is the function $\Chi_t(x) = e^{\frac{2\pi i tx}{n}}$.
We will use $\lambda_G(v)$ to denote the eigenvalue of $v$ which is an eigenvector of the adjacency matrix of graph $G$. For a string $S \in \Sigma^m$, for some alphabet $\Sigma$, and a set $I \subseteq [m]$, we denote its restriction to the set of coordinates in $I$, by $S|_I$.


\subsection{The Graph $C_{n,\ell,\alpha}$ and its Spectrum}

\begin{definition}[Johnson-approximating Cayley Graph $C_{n,\ell,\alpha}$]
Let $\alpha$ be a number in $(0,1)$ and $\ell$ be a positive integer. Let $n$ be a positive integer such that $n > \ell$. 
\begin{enumerate}
\item The vertex set of $C_{n,\ell,\alpha}$ is $[n]^{\ell}$. We will drop the subscript $(n,\ell,\alpha)$ in $C_{n,\ell,\alpha}$ when these parameters are clear from context.
\item The edges are described by the following random process. For a vertex $X = (x_1,\ldots,x_\ell), x_i \in [n]$, the distribution over the neighbours of $X$ is described by: choose $(y_1,\ldots,y_\ell)$ uniformly at random from $[n]^{\ell}$ and $b = (b_1,\ldots,b_{\ell}) \sim \{0,1\}^{\ell}$ such that the Hamming weight of $b$ equals $\alpha \ell$, and output $Z = (x_1 + b_1 \cdot y_1,\ldots,x_l + b_\ell \cdot y_\ell)$ as the neighbour of $X$.
\end{enumerate}
\end{definition}


It is easy to verify that the graph defined above is a weighted Cayley graph with vertex set being the elements of the group $[n]^\ell = (\Z/n\Z)^{\ell}$. The natural group operation associated with this set is component-wise addition modulo $n$. 

We next discuss the spectral properties of $C_{n,\ell\alpha}$. Overloading notations, we let $C = C_{n,\ell,\alpha}$ also denote the normalized adjacency matrix of the graph $C$. Note firstly that the eigenvectors of $C$ are the characters of the group $[n]^\ell$, namely $\raisebox{2pt}{$\chi$}_{T}$ where $T = (T_1,\ldots,T_\ell) \in [n]^{\ell}$ defined as $\Chi_T(x) = \Chi_{T_1}(x_1) \cdot \ldots \cdot \Chi_{T_{\ell}}(x_\ell)$ for $x\in [n]^{\ell}$ 
We next define the degree of an eigenvector.

\begin{definition}[Degree of $\Chi_T$]
For all $T \in [n]^\ell$, where $T = (T_1,\ldots,T_\ell)$, define the degree of $\chi_{T}$ as:
$${\sf deg}(\chi_T) = |T| := |\{i \mid T_i \neq 0\}|.$$
\end{definition}

The following lemma asserts that the eigenvalue of $\chi_T$ is roughly $(1-\alpha)^{|T|}$:
\begin{lemma}\label{lem:eigenval}
Let $\lambda_C(\Chi_T)$ denote the eigenvalue of $C$ corresponding to the eigenvector $\Chi_T$ for $T \in [n]^\ell$. We have that, 
$$
\lambda_{C}(\raisebox{2pt}{$\chi$}_T) = 
\begin{cases}
\frac{\binom{\ell - |T|}{(1 - \alpha)\ell - |T|}}{\binom{\ell}{(1 - \alpha)\ell}}, ~~~~|T| \leq (1 - \alpha)\ell\\
0, ~~~~~~~~~~~~~~~~~~~~~~~\text{otherwise.} 
\end{cases}
$$
\end{lemma}

\begin{proof}
Let $T = (T_1, \ldots, T_\ell)$. For all $X \in [n]^\ell$, we have that,

\begin{align*}
C \cdot \Chi_T(X) = \E_{y,b}[\Chi_T(x_1+b_1 y, \ldots, x_l + b_l y_l)] 
&= \Chi_T(X) \E_{y,b}[\Chi_{T_1,\ldots,T_l}(b_1y_1,\ldots,b_ly_l)] \\
&= \Chi_T(X) \E_{y,b}[\Chi_{(b_1 T_1,\ldots,b_l T_l)}(y)].
\end{align*}

For $y,S \in [n]$ and $S \neq 0$, we know that the eigenvector $\Chi_S$ is orthogonal to the eigenvector $\Chi_0$, equivalently that $\E_y[\Chi_S(y)] = 0$, whereas if $S = 0$ then $\E_y[\Chi_S(y)] = 1$. So we get that,

\begin{align*}
\lambda_{C}(\Chi_T) = \E_{y,b}[\Chi_{(b_1 T_1,\ldots,b_\ell T_\ell)}(y)] 
&= \Pr_b[(b_1T_1,\ldots,b_\ell T_\ell) = 0^{\ell}] \\
&= \begin{cases}
\frac{\binom{l - |T|}{(1 - \alpha)l - |T|}}{\binom{l}{(1 - \alpha)l}}, ~~~~|T| \leq (1 - \alpha)\ell\\
0, ~~~~~~~~~~~~~~~~~~~~~~~\text{otherwise.} 
\end{cases}.
\end{align*}
\end{proof}

\subsection{Analyzing non-expanding sets of the Johnson graph}
Since we want to deal with sets in the Johnson graph $J_{n,\ell,\alpha}$ we will only consider ``permutation-invariant'' sets on $C_{n,\ell,
\alpha}$. Notice that the vertices of the Johnson graph are subsets of $[n]$ of size $\ell$, whereas the vertices of the Johnson-approximating graph are ordered $\ell$-tuples of $[n]$. Therefore, given a set $S$ in the Johnson graph, it has a natural mapping to the set $S'$ which is a subset of the vertices of the Johnson-approximating graph $C$, 
$S' := \{(x_{\pi(1)},\ldots,x_{\pi(\ell)}) \mid \pi: [\ell] \rightarrow [\ell], \{x_1,\ldots,x_{\ell}\} \in S\}$. This leads to the following definition:

\begin{definition}[Permutation-invariance]\label{def:perm-inv}
We say that a set $S \subseteq C_{n,\ell,\alpha}$ is permutation-invariant if for all permutations $\pi \in \mathcal{S}_\ell$ and all $X = (x_1,\ldots,x_\ell) \in S$ we have that $X_{\pi} = (x_{\pi(1)},\ldots,x_{\pi(\ell)})$ belongs to $S$. Similarly, a function $F: V(C) \rightarrow \R$ is permutation invariant if for all inputs $X = (x_1,\ldots,x_{\ell})$, we have that $F(x_1,\ldots,x_\ell) = F(x_{\pi(1)},\ldots,x_{\pi(\ell)})$, for all $\pi\in S_{\ell}$. Further let $\A_{inv}$ denote the set of axioms that $F$ is permutation-invariant, that is, 
\[\A_{inv} := \{F(x_1,\ldots,x_\ell) = F(x_{\pi(1)},\ldots,x_{\pi(\ell)})\}_{\pi \in \mathcal{S}_\ell, X \in [n]^\ell}.\]
\end{definition}

Since the set of vertices in $C$ that correspond to some set of vertices in $J$ are permutation invariant it will be enough to focus are attention on these special sets and from now on whenever we refer to a set in $V(C)$, the reader can assume that it is permutation-invariant.

To analyze non-expanding sets of $C$, we will consider permutation-invariant functions $F: V(C) \rightarrow [0,1]$. Typically one would consider $0/1$-valued functions $F$, where $F$ is the indicator function of a set $S$, i.e. $F(X) = 1$ when $X \in S$. But since we need to analyze ``approximate-sets'' (the indicator function is approximated by a polynomial that takes values close to $0/1$), $F(X)$ could take any value between $[0,1]$.

Recall that the Fourier decomposition of $F$ gives us that, $F(X) = \sum_T \hat{F}(T)\Chi_T(X)$. We will now define the following for a function $F$:

\begin{definition}\label{def:f_i}
Given a function $F\colon V(C)\to [0,1]$, we have the following level decomposition:
\begin{enumerate}
\item We write 
$F = F_0 + \ldots + F_\ell$,
where $F_i(X) = \sum_{T: |T| = i} \hat{F}(T) \Chi_T(X)$. We will call $F$ a level $i$ function, if its Fourier decomposition has degree $i$ characters only, i.e. $\hat{F}(T) = 0,$ for all $T$ such that $|T| \neq i$.
\item Let $f_{i,F}: [n]^i \rightarrow \R$ be a function defined as, 
$$f_{i,F}(x_1,\ldots,x_i) := \sum_{T_1,\ldots,T_i \in ([n] \setminus 0)^i}\hat{F}(T_1,\ldots,T_i,0,\ldots)\chi_{T_1,\ldots,T_i}(x_1,\ldots,x_i),$$ 
\end{enumerate}
\end{definition}



Let $X = (x_1,\ldots,x_{j}) \in [n]^{j}$ and $I$ be a subset of $\{1,\ldots,j\}$. Let $I = \{k_1,\ldots,k_{|I|}\}$ where $k_1 < k_2 < \ldots < k_{|I|}$. We will use $X|_I$ to denote the ordered tuple of elements $(x_{k_1},\ldots,x_{k_{|I|}})$. We will now state some simple properties of $F$ that are implied by permutation-invariance.

\begin{lemma}\label{lem:props}
For all functions $F:[n]^\ell \rightarrow \R$ that are permutation-invariant, we have that:
\begin{enumerate}
\item $\hat{F}(T_1,\ldots, T_\ell) = \hat{F}(T_{\pi(1)},\ldots,T_{\pi(\ell)})$, for all $(T_1,\ldots,T_\ell) \in [n]^\ell$ and all permutations $\pi:[\ell] \rightarrow [\ell]$.

\item The functions $F_i$ and $f_{i,F}$ are also permutation-invariant.

\item $F_i(X) = \sum\limits_{\substack{I \subseteq [\ell] \\ |I| = i}} f_{i,F}(X|_{I})$.
\end{enumerate}
\end{lemma}
\begin{proof}
A straightforward manipulation of the definitions.
\end{proof}

We note that the Fourier coefficients $\hat{F}(T), \forall T \in \C$, the level-$i$ functions $F_i(X)$ and $f_i(X)$ are all linear functions of the indeterminates $\{F(X)\}$,
hence we shall also think of them as indeterminates 
when arguing about \sos proofs.

\subsection{Restrictions}
In this section, we define restrictions of functions $F$ and state several lemmas related to them.
\begin{definition}[$r$-restricted subcubes of $C$]
Given an ordered tuple, $A = (a_1,\ldots,a_r)$ for $a_i \in [n]$ and $r \leq \ell-1$, we let $C|_{A}$ denote the subset of vertices of $C$ whose first $r$ coordinates are restricted to be $(a_1,\ldots,a_r)$. We call such a subset an $r$-restricted subcube of $C$.
\end{definition}


\begin{definition}[Restrictions]
Given a function $F:[n]^\ell \rightarrow \R$ and an ordered tuple, $A = (a_1,\ldots,a_r)$ for $a_i \in [n]$ and $1 \leq r \leq \ell-1$, we define the restricted function $F|_A: [n]^{\ell - r} \rightarrow \R$ as, 
$$F|_A(x_1,\ldots,x_{\ell-r}) = F(a_1,\ldots,a_r, x_1,\ldots,x_{\ell-r}).$$
Further, let $\delta_A(F)$ denote the mass of the function restricted to $A$, that is,
$$\delta_A(F) := \delta(F|_A) = \E\limits_{X \in [n]^{\ell - r}}[F|_A(X)].$$
For convenience, when $A = \phi$ ($r = 0$), define $F|_A(X) := F(X)$, and $\delta_A(F) := \delta(F) = \E_{X \in [n]^\ell}[F(X)]$.
\end{definition}



The following lemma gives a relation between the level functions of a function $F$ and the level function of its restrictions.
\begin{lemma}\label{lem:restriction}
Let $F$ be a permutation-invariant function on $V(C)$. Then we have the following:
\begin{enumerate}
\item For all $a \in [n]$ and for all $i$ such that $0 \leq i \leq \ell-1$, and all $X \in [n]^i$, we have that,
\[f_{i+1,F}(a, X) = f_{i,F|_{\{a\}}}(X) - f_{i,F}(X).\]

\item For all integers $i$ such that $0 \leq i \leq \ell$ and for all $X \in [n]^i$, we get an inclusion-exclusion formula for $f_i(X)$ in terms of restrictions of $F$: 
\[f_{i,F}(X) = \sum_{B \subseteq \{1,\ldots,i\}} (-1)^{i - |B|} \delta_{X|_B}(F),\]
where $X|_B$ is the ordered tuple of elements of $X$ restricted to the indices in $B$.
\end{enumerate}
\end{lemma}

\begin{proof}
We prove each item separately.
\paragraph{Proof of the first item.}
Using the definition, we can expand out $f_{i+1,F}$ to get that,

\[f_{i,F}(a, X) = \sum_{(T_1,\ldots,T_{i+1}) \in ([n] \setminus 0)^{i+1}}\hat{F}(T_1,\ldots,T_{i+1},0,\ldots, 0)\Chi_{T_1,\ldots,T_{i+1}}(a, X).\]

We can split this sum into two parts, one where $T_{1}$ can take \emph{any} value (even $0$) and the 
second where $T_{1} = 0$. We get that,

\begin{align*}
f_{i,F}(a,X) &= \sum_{ \substack{T_1 \in [n] \\ T \in ([n] \setminus 0)^i}}\hat{F}(T_1,T, 0,\ldots,0,\ldots, 0)\Chi_{T_1,T}(a,X) 
- \sum_{ \substack{T_1 = 0 \\ T \in ([n] \setminus 0)^i}}\hat{F}(0, T ,0,\ldots, 0)\Chi_{T}(X).
\end{align*}
We will show that the first term equals $f_{i, F|_{\{a\}}}(X)$ and the second term equals $f_{i,F}(X)$. This implies the conclusion needed.

For the first term we have that, 
\begin{align}
&\sum_{ \substack{T_1 \in [n] \\ T \in ([n] \setminus 0)^i}}\hat{F}(T_1, T,0,\ldots,0)\Chi_{T_1,T}(a, X) \nonumber \\
= &\sum_{ \substack{T_1 \in [n] \\ T \in ([n] \setminus 0)^i}} \E_{\substack{Y_1 \in [n], \\ Y \in [n]^{\ell - 1}}}\left[F(Y_1,Y)\overline{\Chi_{T_1}}(Y_1)\overline{\Chi_{(T,0,\ldots,0)}(Y)}\right]\Chi_{T_1}(a)\Chi_{T}(X) \nonumber \\
= &\sum_{T \in ([n] \setminus 0)^i} \E_{\substack{Y_1 \in [n], \\ Y \in [n]^{\ell - 1}}}\left[F(Y_1,Y)\overline{\Chi_{(T,0,\ldots,0)}(Y)}\sum_{T_1 \in [n]}\Chi_{T_1}(a-Y_1)\right]\Chi_{T}(X). \label{eq:last}
\end{align}

We now have that $\sum_{T_1 \in [n]}\Chi_{T_1}(a-Y_1) = 0$ if $Y_1 \neq a$ and equals $n$ otherwise. Using this fact we get that~\eqref{eq:last} is equal to
\begin{align*}
\sum_{T \in ([n] \setminus 0)^i} \frac{1}{n} \cdot \E_{Y \in [n]^{\ell - 1}}\left[F(a,Y)\overline{\Chi_{(T,0,\ldots,0)}(Y)}\cdot n\right]\Chi_{T}(X) 
= \sum_{T \in ([n] \setminus 0)^i} \widehat{F|_{\{a\}}}(T,0,\ldots,0)\Chi_{T}(X),
\end{align*}
which is equal to $f_{i, F|_{\{a\}}}(X)$ by definition.


As for the second term, it is equal to
\begin{align*}
&\sum_{ \substack{T_1 = 0 \\ (T_2,\ldots,T_{i+1}) \in ([n] \setminus 0)^i}}\widehat{F}(0, T_2,\ldots,T_{i+1},0,\ldots, 0)\Chi_{T_2,\ldots,T_{i+1}}(X)\\
= &\sum_{(T_2,\ldots,T_{i+1}) \in ([n] \setminus 0)^i}\hat{F}(T_2, \ldots,T_{i+1},0,\ldots, 0)\Chi_{T_2,\ldots,T_{i+1}}(X),
\end{align*}
since by Lemma~\ref{lem:props} (1) we have that $\hat{F}(0, T_2,\ldots,T_{i+1}, 0, \ldots) = \hat{F}(0, T_2,\ldots,T_{i+1}, 0, \ldots)$. Since the last equality is the definition of $f_{i,F}(X)$, the conclusion follows.



\paragraph{Proof of the second item.}
We will prove this claim by induction on $i$. For the base case of $i = 0$, by definition, we have that,

\[f_{0,F}(\phi) = \hat{F}(0,\ldots,0) = \E_{X \in [n]^\ell}[F(X)] = \delta(F) = \delta_{\phi}(F) = \sum_{B \subseteq \phi} \delta_{\phi|_B}(F).\]

Now let us assume that for all permutation-invariant functions $G$ the claim holds for $i-1$, i.e. for all $X \in [n^{i-1}]$, we have that $f_{i-1,G}(X) = \sum_{B \in \{1,\ldots,i-1\}} (-1)^{i-1-|B|} \delta_{X|_{B}}(G)$. Now we will prove the  claim for $f_{i,F}$, thus completing the induction.
Let $X = (x_1, X')$, where $X \in [n]^i, x_1 \in [n]$ and $X' \in [n]^{i-1}$. Then by property (1) of the same lemma, we have that,
\[f_{i,F}(X) = f_{i,F|_{\{x_1\}}}(X') - f_{i,F}(X').\]

Expanding the RHS using the induction hypothesis on the functions $F|_{\{x_1\}}$ and $F$, we get that,

\begin{align*}
f_{i,F}(X) &= \sum_{B' \in \{1,\ldots,i-1\}} (-1)^{i-1-|B'|} \delta_{X'|_{B'}}(F|_{\{x_1\}}) - \sum_{B' \in \{1,\ldots,i-1\}} (-1)^{i-1-|B'|} \delta_{X'|_{B'}}(F) \\
&= \sum_{B' \in \{1,\ldots,i-1\}} (-1)^{i-(1+|B'|)} \delta_{(x_1,X'|_{B'})}(F) + \sum_{B' \in \{1,\ldots,i-1\}} (-1)^{i-|B'|} \delta_{X'|_{B'}}(F) \\
&= \sum_{B \in \{1,\ldots,i\}: 1 \in B} (-1)^{i-|B|} \delta_{X|_B}(F) + \sum_{B \in \{1,\ldots,i\}: 1 \notin B} (-1)^{i-|B|} \delta_{X|_{B}}(F) \\
&= \sum_{B \in \{1,\ldots, i\}} (-1)^{i-|B|} \delta_{X|_B}(F).
\end{align*}

This completes the inductive step and the proof of the lemma.
\end{proof}



\subsection{Calculating $2$-nd Moments}
Let $W_i(F)$ denote the weight of $F$ on the $i^{th}$ level, that is, let $\eta_i = W_i[F] = \E_X[F_i(X)^2]$. Let $\delta(F)$ denote $\E_X[F(X)]$. Now we will derive some lemmas about the second moments. We get the following relation between the level-$i$ Fourier weight $\eta_i$ and the $\ell_2$-norm of the $f_i$'s.

\begin{lemma}\label{lem:3.2}
For all $i \in [\ell]$, given the variables $\{F(X)\}_{X \in V(C)}$:
$$\cA_{inv} ~~\vdash_2 ~~\E_{X \in [n]^{i}}[f_{i,F}(X)^2] =  \frac{\eta_{i}}{\binom{\ell}{i}}.$$
\end{lemma}

\begin{proof}
Since $F_i(A) = \sum_{I \subseteq [\ell]}f_i(A|_I)$, we have that, 
\begin{align*}
    \eta_i=\E_{A \sim [n]^{\ell}}[F_i(A)^2] = 
    \E_A[(\sum_I f_i(A_I))^2] 
    &= \sum_I \E_A[ f_i(A_I)^2] + \sum_{I \neq I'} \E_{A}[f_i(A_I)f_i(A_{I'})] \\
    &= \sum_I \E_A[ f_i(A_I)^2] + 0 \\
    &= \E_{(x_1,\ldots,x_i)}[f_i(X)^2] \cdot \binom{\ell}{i},
\end{align*}
and rearranging implies the lemma. 
\end{proof}




\begin{lemma}\label{lem:p.r.-weight}
For all $i \in [\ell]$, for all $a,b \in [i]$, for all $A \in (\{0,1\}^k)^a$, given the indeterminates $\{F(X)\}_{X \in V(C)}$:
\begin{align*}
\cA_{inv}(F) ~~~\vdash_2 ~~~ &\E_{X \in (\{0,1\}^k)^{i-a}}[f_i(A,X)^2] \leq \frac{\delta(F^{ 2}|_A)}{{l-a \choose i-a}}.
\end{align*}
\end{lemma}

\begin{proof}
By Lemma~\ref{lem:3.2}, we have that,
\begin{align}
\E_{X \in (\{0,1\}^k)^{i-a}}[f_i(A,X)^2] = \E_{X \sim (\{0,1\}^k)^{i-a}}[f_{i-a,F|_A}(B)^2] = \frac{W^{i-a}(F|_A)}{{l-a \choose i-a}}. \label{5.5:eq1}
\end{align}

We also have that,

\begin{align}
W^{i-a}(F|_A) \leq \E_{X \in (\{0,1\}^k)^{i-a}}[F|_A(X)^2] = \delta(F^{ 2}|_A), \label{5.5:eq2}
\end{align}
where the first inequality is a degree 2 SoS-inequality. The proof is concluded by combining~\eqref{5.5:eq1},~\eqref{5.5:eq2}.
\end{proof}


\subsection{Bounding $4$-th moments}
The next part of the proof is to prove a lower 
bound on the fourth moment of $F_i$ (using the fact
it is correlated with a Boolean valued function, 
hence is supposed to have a high $4$th moment) 
as well as an upper bound on the fourth moment of 
$F_i$ (using the fact $F$ is pseudo-random). We 
begin with the following lemma which gives us the lower bound:
\begin{lemma}[Lower Bound]\label{lem:lower-bound}
For all real $\epsilon > 0$ given the indeterminates $\{F(X)\}_{X \in V(C)}$ we have that,
$$\E_X[F_i^4(X)] \geq 4\epsilon^3 \eta_i - 3\epsilon^4\delta(F) + B(F),$$
where $B(F) = 4\eps^3\E_X[(F^{ 3} - F)F_i]+3\eps^4 \E_X[F - F^{ 4}]$.
\end{lemma}

\begin{proof}
We have that, 
\begin{align}\label{eq:eta-i}
\eta_i = \ip{F,F_i} = \E_X[F(X)^3 \cdot F_i(X)] + \E_X[(F(X) - F(X)^3) \cdot F_i(X)],
\end{align}
since $F = \sum F_i$ is an orthogonal decomposition of $F$. Using Fact~\ref{fact:sos-hol} on the first term we get:
\begin{align*}
\E_X[F(X)^3 \cdot F_i(X)] &= 
\E_X[F(X)^3 \cdot F_i(X)]\\
&\leq  
\frac{3\epsilon}{4}\E_X[F(X)^4] + \frac{1}{4\epsilon^3}\E_X[F_i(X)^4],\\
&= \frac{3\epsilon}{4}\E_X[F(X)] +\frac{3\epsilon}{4}\E_X[F(X)^4 - F(X)] +  \frac{1}{4\epsilon^3}\E_X[F_i(X)^4],
\end{align*}
for all real $\epsilon > 0$. Plugging this into~\eqref{eq:eta-i} and rearranging gives the lemma.
\end{proof}


We will now prove an upper bound on $\E_X[F_i^4(X)]$. We have the following lemma which bounds the fourth moment of $F_i$ in terms of the second moment of $F_i$:

% \begin{lemma}
% Let $i \leq r$ and $F$ be the indicator of a basis invariant set of $J_{k,l,\alpha}$, with $\eta = \E_X[F_i^2(X)]$ and $\delta = \E_X[F(X)]$. Then we have that,
% $$\E_X[F_i^4(X)] \leq \exp(i)\gamma\eta_i.$$
% \end{lemma}


\begin{lemma}[Upper Bound]\label{lem:upper-bound}
For all integers $i \in [\ell]$, there exist positive constants $c_1, \ldots, c_i \leq \exp(i)$ such that for all $F$:
\[\cA_{inv} \,\, \vdash_{O(1)}\,\, \E_X[F_i(X)^4] \leq \exp(i)\gamma \eta_i + \sum_{j = 0}^i c_j \ell^{i}\E_{a \sim [n]^{j}}[\E_{b \sim [n]^{\ell-j}}[f_i(a,b)^2]\cdot (\delta(F^{ 2}|_a) - \gamma)].\] 
\end{lemma}


\begin{proof}
When the set $X$ and the function $F$ is clear from context, we use $f_i(I)$ in place of $f_i(X|_I)$. Thus,
\begin{align*}
\E_X[F_i^4(X)] &= \sum_{\substack{ I_1,\ldots,I_4 \subseteq [l] \\ |I_j| = i, \forall j\in [4]}}\E_X[ f_i(I_1) f_i(I_2) f_i(I_3) f_i(I_4)].   
\end{align*}
Denote, $$D(i,d) = \{(I_1,\ldots,I_4) \mid \forall j \in [4], I_j \subseteq [l], |I_j| = i, |I_1 \cup I_2 \cup I_3 \cup I_4| = d\}.$$
Note that for all $d \in [i,4i]$,
\[|D(i,d)| \leq {\ell \choose d}\cdot {d \choose i}^4 \leq \left(\frac{\ell e}{d}\right)^d \left(\frac{de}{i} \right)^{4i} \leq \ell^d \exp(i).\]

We will prove the following lemma about the four-wise products, for $(I_1,\ldots,I_4) \in D(i,d)$. We defer
the proof to Section~\ref{sec:pf_lem-prod}.
\begin{lemma}\label{lem-prod}
Let $i,d$ be integers such that $i \leq d \leq 4i$ and let $I_1,\ldots,I_4 \subseteq [l]$, such that $(I_1,I_2,I_3,I_4) \in D(i,d)$. Then there exist integers $t_1,t_2 \in [i]$ such that for all $\gamma$:
\begin{align*}
&\E_X[f_i(I_1)f_i(I_2)f_i(I_3)f_i(I_4)] \\
&\leq \left(\frac{O(i)}{\ell}\right)^{d}\gamma\eta_i + \left(\frac{O(i)}{\ell}\right)^{d-i}\E_{a_1 \sim [n]^{t_1}}[\E_{b_1 \sim [n]^{i-t_1}}[f_i(a_1,b_1)^2]\cdot (\delta(F^{ 2}|_{a_1}) - \gamma)]  \\ 
&+ \left(\frac{O(i)}{\ell}\right)^{d-i}\E_{a_2 \sim [n]^{t_2}}[\E_{b_2 \sim [n]^{i-t_2}}[f_i(a_2,b_2)^2](\delta(F^{ 2}|_{a_2}) - \gamma)]. 
\end{align*}
\end{lemma}
We now use Lemma~\ref{lem-prod} to finish the proof of Lemma~\ref{lem:upper-bound}. Indeed, we get that
\begin{align*}
&\E_X[F_i^4(X)] \\
&= \sum_{\substack{ I_1,\ldots,I_4 \subseteq [l] \\ |I_j| = i, \forall j\in [4]}}\E_X[ f_i(I_1) f_i(I_2) f_i(I_3) f_i(I_4)] \\
&\leq \sum_{d = i}^{4i} \sum_{\substack{ I_1,\ldots,I_4 \subseteq [l] \\ (I_1,\ldots,I_4) \in D(i,d)}} \E_X[ f_i(I_1) f_i(I_2) f_i(I_3) f_i(I_4)] \\
&\leq \sum_{d = i}^{4i} \sum_{\substack{ I_1,\ldots,I_4 \subseteq [l] \\ (I_1,\ldots,I_4) \in D(i,d)}} \left(\frac{O(i)}{\ell}\right)^{d}\gamma\eta_i +  \left(\frac{O(i)}{\ell}\right)^{d-i}\E_{a_1 \sim [n]^{t_1}}[\E_{b_1 \sim [n]^{i-t_1}}[f_i(a_1,b_1)^2]\cdot (\delta(F^{ 2}|_{a_1}) - \gamma)]  \\ 
& ~~~~~~~~~~~~~~~~~~~~~~~~~~~~~~~~~+\left(\frac{O(i)}{\ell}\right)^{d-i}\E_{a_2 \sim [n]^{t_2}}[\E_{b_2 \sim [n]^{i-t_2}}[f_i(a_2,b_2)^2](\delta(F^{ 2}|_{a_2}) - \gamma)] \\
&\leq \exp(i)\gamma \eta_i + \sum_{j = 0}^i c_j \ell^{i}\E_{a \sim [n]^{j}}[\E_{b \sim [n]^{i - j}}[f_i(a,b)^2]\cdot (\delta(F^{ 2}|_a) - \gamma)],  
\end{align*}
for some positive constants $c_j \leq \exp(i)$, where we have used that $|D(i,d)| \leq \ell^d \exp(i)$.
\end{proof}

\subsection{Proof of Lemma~\ref{lem-prod}}\label{sec:pf_lem-prod}
We begin by analyzing the simple case in which there is an index that appears in only one of the $I_j$'s, 
and show that in this case the expectation is just equal to $0$:
\begin{proposition}
If there is an index $j$ that appears in only one of the $I_k$s, then,
$$\E_X[ f_i(I_1) f_i(I_2) f_i(I_3) f_i(I_4)] = 0.$$
\end{proposition}

\begin{proof}
Without loss of generality we can assume that $1$ appears only in $I_4$. We have that,
\begin{align*}
\E_X[ f_i(I_1) f_i(I_2) f_i(I_3) f_i(I_4)] &= \E_{x_1,\ldots,x_l}[ f_i(I_1) f_i(I_2) f_i(I_3) f_i(I_4)] \\
&= \E_{x_2,\ldots,x_l}[f_i(I_1) f_i(I_2) f_i(I_3) \E_{x_1}[f_i(I_4) \mid x_2 = x_2',\ldots, x_l = x_l']]\\
&= \E_{x_2,\ldots,x_l}[f_i(I_1) f_i(I_2) f_i(I_3) \delta_{(x_2',\ldots,x_l') \cap I_4}(f_i)] \\
&= 0,
\end{align*}
where the last equality follows by Lemma~\ref{lem:props}. 
%\mitali{point (4): check/add}
%\dor{what is this? resolve?}
\end{proof}
For the rest of the proof we assume each index in $I_1\cup I_2\cup I_3\cup I_4$ appears in at least two sets. Without loss of generality we can assume that $\cup_{j = 1}^4 I_j = [d]$. Let $H_2, H_3, H_4 \subseteq [d]$ be the set of elements that appear in $2,3,4$ of the sets $(I_1,\ldots,I_4)$ respectively. Furthermore let $H_2 = X_{12} \cup X_{13} \ldots X_{34},$ where $X_{jk}$ is the set of elements that occur in $I_j$ and in $I_k$. We will abuse notation and use $H_j$ to also denote the set of random variables $(x_{j_1},\ldots, x_{j_m}), m \in H_j$ and $x_m \sim \{0,1\}^k$. Firstly we get that, 

\begin{align}
&\E_{H_2}[ f_i(I_1) f_i(I_2) f_i(I_3) f_i(I_4)] \nonumber\\
= &\E_{X_{12},\ldots,X_{34}}[f_i(I_1) f_i(I_2) f_i(I_3) f_i(I_4)] \nonumber\\
= &\E_{\substack{X_{13},X_{14},\\ X_{23},X_{24}}}[\E_{X_{12}}[f_i(I_1) f_i(I_2)]\E_{X_{34}}[f_i(I_3) f_i(I_4)] ] \nonumber\\
\leq &\frac{\eps}{2}\E_{\substack{X_{13},X_{14},\\ X_{23},X_{24}}}[\E_{X_{12}}[f_i(I_1) f_i(I_2)]^2] +  \frac{1}{2\eps}\E_{\substack{X_{13},X_{14},\\ X_{23},X_{24}}}[\E_{X_{34}}[f_i(I_3) f_i(I_4)]^2] \label{CS1} \\
\leq &\frac{\eps}{2}\E_{\substack{X_{13},X_{14},\\ X_{23},X_{24}}}[\E_{X_{12}}[f_i(I_1)^2]\E_{X_{12}}[f_i(I_2)^2]] +  \frac{1}{2\eps}\E_{\substack{X_{13},X_{14},\\ X_{23},X_{24}}}[\E_{X_{34}}[f_i(I_3)^2]\E_{X_{34}}[f_i(I_4)^2]]\label{CS2} \\
\leq &\frac{\eps}{2}\E_{X_{12},X_{13},X_{14}}[f_i(I_1)^2]\E_{X_{12},X_{23},X_{24}}[f_i(I_2)^2] +  \frac{1}{2\eps}\E_{X_{34},X_{13},X_{23}}[f_i(I_3)^2]\E_{X_{34},X_{14},X_{24}}[f_i(I_4)^2] \nonumber \\
= &\frac{\eps}{2}\E_{H_2}[f_i(I_1)^2]\E_{H_2}[f_i(I_2)^2] +  \frac{1}{2\eps}\E_{H_2}[f_i(I_3)^2]\E_{H_2}[f_i(I_4)^2],\label{last-eq} 
%\leq &\sqrt{\T[\E_{H_2}[f_i(I_1)^2]\E_{H_2}[f_i(I_2)^2]]} \cdot \sqrt{\T[\E_{H_2}[f_i(I_3)^2]\E_{H_2}[f_i(I_4)^2]]}
\end{align}
where inequalities~\eqref{CS1},~\eqref{CS2} follow by Cauchy Schwarz (Lemma~\ref{CS1-prelim} and Lemma~\ref{CS2-prelim}). 

%and the last inequality follows by setting $\epsilon = \sqrt{\frac{\T[\E_{H_2}[f_i(I_3)^2]\E_{H_2}[f_i(I_4)^2]]}{\T[\E_{H_2}[f_i(I_1)^2]\E_{H_2}[f_i(I_2)^2]]}}$. 

\bigskip


We will now bound the expression $\E_X[f_i(I_1)f_i(I_2)f_i(I_3)f_i(I_4)]$. We have that,

\begin{align}
&\E_X[f_i(I_1)f_i(I_2)f_i(I_3)f_i(I_4)] \nonumber \\
= &\E_{H_4,H_3} [\E_{H_2}[f_i(I_1)f_i(I_2)f_i(I_3)f_i(I_4)]]\nonumber \\
\leq &\frac{\eps}{2}\E_{H_4,H_3}[\E_{H_2}[f_i(I_1)^2]\E_{H_2}[f_i(I_2)^2]] +  \frac{1}{2\eps}\E_{H_4,H_3}[\E_{H_2}[f_i(I_3)^2]\E_{H_2}[f_i(I_4)^2]]\nonumber\\
= &\frac{\eps}{2}\E_{A}[\E_{B_1}[f_i(I_1)^2]\E_{B_2}[f_i(I_2)^2]] +  \frac{1}{2\eps}\E_{C}[\E_{D_1}[f_i(I_3)^2]\E_{D_2}[f_i(I_4)^2]],\nonumber\\
= &\frac{\eps}{2}\E_{A}[\E_{B_1}[f_i(a, b_1)^2]\E_{B_2}[f_i(a, b_2)^2]] +  \frac{1}{2\eps}\E_{C}[\E_{D_1}[f_i(c,d_1)^2]\E_{D_2}[f_i(c,d_2)^2]],\label{eq:hypcont}
\end{align}
with $A = H_4 \cup (H_3 \cap I_1 \cap I_2)$, 
$B_1 = (H_2 \cap I_1) \cup ((H_3 \cap I_1) \setminus I_2)$,
$B_2 = (H_2 \cap I_2) \cup ((H_3 \cap I_2) \setminus I_1)$, $C = H_4 \cup (H_3 \cap I_3 \cap I_4)$, $D_1 = (H_2 \cap I_3) \cup ((H_3 \cap I_3) \setminus I_4)$,
$D_2 = (H_2 \cap I_4) \cup ((H_3 \cap I_4) \setminus I_3)$. The sets $a,b_1,b_2,c,d_1,d_2$ are instantiations of the corresponding random variables, $A, \ldots,D_2$. We have that,

\begin{align*}
\E_{A}[\E_{B_1}[f_i(a, b_1)^2]\E_{B_2}[f_i(a, b_2)^2]] &\leq \E_{A}[\frac{\delta(F^{ 2}|_a)}{{\ell-|A| \choose i - |A|}}\E_{B_2}[f_i(a,b_2)^2]] \\
&= \frac{1}{{\ell-|A| \choose i - |A|}}\E_{A}[\gamma \E_{B_2}[f_i(a,b_2)^2]] + \E_A[(\delta(F^{ 2}|_A) - \gamma)\left(\frac{\E_{B_2}[f_i(a,b_2)^2]}{{\ell-|A| \choose i - |A|}}\right)],\\
&\leq \frac{\gamma}{{\ell-|A| \choose i - |A|}}\cdot \frac{\eta_i}{{\ell \choose i}} + \E_A[(\delta(F^{2}|_A) - \gamma)\left(\frac{\E_{B_2}[f_i(a,b_2)^2]}{{\ell-|A| \choose i - |A|}}\right)],
\end{align*}
where in the first inequality we have used Lemma~\ref{lem:p.r.-weight} to bound $\E_{B_1}[f_i(a,b_1)^2]$ for each $a \sim A$ and the last inequality uses Lemma~\ref{lem:3.2}. Similarly we can bound the second term in equation~\eqref{eq:hypcont}:
\[\E_{C}[\E_{D_1}[f_i(c,d_1)^2]\E_{D_2}[f_i(c,d_2)^2]] \leq \frac{\gamma}{{\ell-|C| \choose i - |C|}}\cdot \frac{\eta_i}{{\ell \choose i}} + \E_C[(\delta(F^{2}|_C) - \gamma)\left(\frac{\E_{D_2}[f_i(c,d_2)^2]}{{\ell-|C| \choose i - |C|}}\right)].\]

Plugging in these bounds into equation~\eqref{eq:hypcont} and setting $\eps = \left(\frac{O(i)}{\ell}\right)^{\frac{|A|-|C|}{2}}$ we get:
\begin{align*}
&\E_X[f_i(I_1)f_i(I_2)f_i(I_3)f_i(I_4)] \\
&\leq \left(\frac{O(i)}{\ell}\right)^{\frac{4i - |A|-|C|}{2}}\gamma\eta_i + \E_A[ \left(\frac{O(i)}{\ell}\right)^{\frac{2i-|A|-|C|}{2}}\E_{B_2}[f_i(a,b_2)^2]\cdot (\delta(F^{ 2}|_a) - \gamma)]  \\ 
&+\E_C[\left(\frac{O(i)}{\ell}\right)^{\frac{2i-|A|-|C|}{2}}\E_{D_2}[f_i(c,d_2)^2](\delta(F^{ 2}|_c) - \gamma)]. 
\end{align*}

We know that $4i - |A|-|C| = 2d$ and substituting this gives the lemma statement.




\subsection{Combining the Upper and Lower Bounds on $\|F_i\|_4$: the Level $i$ Inequality}
Combining the upper and lower bounds on $\E[F_i(X)^4]$ we can now prove an upper bound on $\eta_i$.

\begin{lemma}[Level $i$ inequality]\label{lem:leveli-j}
There exist positive constants $c_0,\ldots,c_i \leq \exp(i)$ such that given the indeterminates $\{F(X)\}_{X \in V(C)}$,
\begin{align*}
\cA_{inv} \,\, \vdash_{O(1)}\,\, 
\eta_i \leq \exp(i)\gamma^{1/3}\delta + B(F) + \sum_{j = 0}^i  \frac{c_j\ell^{i}}{\gamma}\E_{a \sim [n]^{j}}[\E_{b \sim [n]^{i-j}}[f_i(a,b)^2]\cdot (\delta(F^{ 2}|_a) - \gamma)],
\end{align*}
where $B(F) =  \frac{4}{3}\E_X[(F - F^{ 3})F_i] + \exp(i)\gamma^{1/3}\E_X[F^{ 4} - F]$.
\end{lemma}

\begin{proof}
Let $\delta = \E_X[F(X)]$ and $\eta_i = \E_X[F_i^2(X)]$. Under the Booleanity axioms on $F$, we have proved an upper and lower bound on ${\E_X[F_i^4(X)]}$,
\begin{itemize}
    \item Upper Bound in Lemma~\ref{lem:upper-bound}:
\begin{equation}\label{eq:ub}
\E_X[F_i(X)^4] \leq \exp(i)\gamma \eta_i + \sum_{j = 0}^i c'_j \ell^{i}\E_{a \sim [n]^{j}}[\E_{b \sim [n]^{i-j}}[f_i(a,b)^2]\cdot (\delta(F^{2}|_a) - \gamma)],
\end{equation}
for some positive constants $c'_j \leq \exp(i)$.

\item Lower Bound in Lemma~\ref{lem:lower-bound}: For all $\eps \in \R$ we get,
\begin{equation}\label{eq:lb}
\E_X[F_i^4(X)] \geq 4\eps^3 \eta_i - 3\eps^4\delta + B_1(F),
\end{equation}
where $B_1(F) = 4\eps^3\E_X[(F^{ 3} - F)F_i]+3\eps^4 \E_X[F - F^{ 4}]$.
\end{itemize}

Combining~\eqref{eq:ub} and~\eqref{eq:lb}, 
setting $\eps = (\exp(i)\gamma)^{1/3}$ and dividing 
by $3\eps^3$ yields
\begin{align*}
\eta_i \leq \exp(i)\gamma^{1/3}\delta + \sum_{j = 0}^i \frac{c'_j}{3\gamma\exp(i)} \ell^{i}\E_{a \sim [n]^{j}}[\E_{b \sim [n]^{i-j}}[f_i(a,b)^2]\cdot (\delta(F^{ 2}|_a) - \gamma)] - \frac{B_1(F)}{3\gamma \exp(i)}.
\end{align*}
The proof is thus completed by setting $c_j = \frac{c'_j}{3\exp(i)}$ and 
$B(F) = \frac{4}{3}\E_X[(F - F^{ 3})F_i] + \exp(i)\gamma^{1/3}\E_X[F^{ 4} - F]$.
\end{proof}


\subsection{Expansion of Pseudorandom sets}
Lemma~\ref{lem:leveli-j} directly implies an expansion theorem for $C_{n,\ell,\alpha}$, as follows:
\begin{theorem}[Expansion Theorem for $C_{n,\ell,\alpha}$]\label{thm:structure-johnson-cayley}
For all $\alpha \in (0,1)$, all integers $\ell \geq 1/\alpha$ and $n \geq \ell$, the following holds: Let $C_{n,\ell,\alpha}$ be the Johnson-approximating Cayley graph. For every constant $\gamma \in (0,1)$ and positive integer $r \leq O(\ell)$, every permutation-invariant function $F: V(C) \rightarrow \R$ that is $(r,\gamma)$-pseudorandom:
\begin{enumerate}
\item $\cA_{inv}(F) \cup \{F(X) \in [0,1]\}_{X \in V(C)} \, \vdash_{O(1)}\, \\
\ip{F, LF} \geq \E[F](1 - (1-\alpha)^{r+1})(1- \gamma^{1/3}\exp(r)) - \sum_{j = 0}^r \frac{ c_j\ell^{j}}{\gamma}\E_{a \sim [n]^{j}}[q_a(F)({\delta(F^{ 2}|_a) - \gamma)}] + B(F).$
\item $\{F(X) \in [0,1]\} \,\,\, \vdash_{2} \,\,\, 0 \leq q_a(F) \leq 1,$
\end{enumerate}
where for all $j \leq r$, $c_j$'s are positive constants $\leq \exp(r)$, for all $a \in [n]^j, j \leq r$, $q_a(F)$ are degree $2$  polynomials and $B(F) = \frac{4}{3}\E_X[(F^{ 3} - F)\sum_{i = 0}^r F_i] + c\E_X[F - F^{ 4}]$, for a positive constant $c \leq \exp(r)\gamma^{1/3}$.
\end{theorem}


\begin{proof}
%\ip{F,LF} = \frac{|E(S,\overline{S})|}{2|E|}
%add axiom F(X) \in [0,1]
We first have that: $\ip{F, LF} = \E[F^{ 2}]  - \ip{F,CF}$, where $L$ and $C$ denote the normalized Laplacian and adjacency matrix of the graph $C_{k,\ell,\alpha}$ respectively. Let $1=\lambda_1 \geq \ldots \lambda_\ell \geq -1$ denote the eigenvalues of $C$. We will henceforth bound $\ip{F,CF}$:
\begin{align*}
\ip{F,CF} = \sum_{i \in [l]} \lambda_i \eta_i 
&\leq \sum_{i = 0}^r \lambda_i\eta_i + (\E[F^{ 2}] - \sum_{i=1}^r \eta_i)\lambda_{r+1} \\
&\leq (1 - \lambda_{r+1}) \sum_{i=0}^r \eta_i + \E[F]\lambda_{r+1},\\
\end{align*}
where in the last step we use that $F(X) \in [0,1]$. Let $\delta$ denote $\E[F]$. Using Lemma~\ref{lem:leveli-j} we get: 
\begin{align}
&\ip{F,CF} \notag\\
&\leq (1 - \lambda_{r+1})\sum_{i=0}^r \left(\exp(i)\gamma^{1/3}\delta + B_{i}(F)+ \sum_{j = 0}^i  \frac{c'_{ij}\ell^{i}}{\gamma}\E_{a \sim [n]^{j}}[\E_{b \sim [n]^{i-j}}[f_i(a,b)^2]\cdot (\delta(F^{2}|_a) - \gamma)] \right) + \delta\lambda_{r+1} \notag\\
&\leq (1 - \lambda_{r+1})(\exp(r)\gamma^{1/3}\delta) + (1-\lambda_{r+1})\sum_i B_i(F) \notag\\+ &\sum_{j = 0}^r  \E_{a \sim [n]^{j}}[ (\delta(F^{ 2}|_a) - \gamma)(\sum_{i = j}^r (1 - \lambda_{r+1})\frac{c'_{ij} \ell^{i}}{\gamma} \E_{b \sim [n]^{i-j}}[f_i(a,b)^2])] +  \delta\lambda_{r+1}, \notag
\end{align}
for some positive constants $c'_{ij} \leq \exp(i)$ and $B_i(F) =  \frac{4}{3}\E_X[(F - F^{ 3})F_i] + \exp(i)\gamma^{1/3}\E_X[F^{ 4} - F]$.

Let us now look at the multipliers of $\delta(F^{2}|_a) - \gamma$, for each $a \in [n]^j$ consider the polynomial: $q_a(F) = \frac{1}{c_j\ell^j}\sum_{i = j}^r (1 - \lambda_{r+1})c'_{ij} \ell^{i} \E_{b \sim [n]^{i-j}}[f_i(a,b)^2]$, for fixed constants $c_j \leq \exp(r)$ to be determined later. Using Lemma~\ref{lem:p.r.-weight} to bound $\E_{b \sim [n]^{i-j}}[f_i(a,b)^2]$ we get that:
\begin{align*}
\{F(X) \in [0,1]\} \vdash_2 ~~~~0 \leq q_a(F) \leq \frac{1}{c_j}\sum_{i = j}^r (1 - \lambda_{r+1})c'_{ij}\ell^{i-j} \frac{ \delta(F^{2}|_a)}{{\ell - j \choose i - j}} \leq 1,
\end{align*}
for $c_j = \sum_{i = j}^r (1 - \lambda_{r+1})c'_{ij}\frac{\ell^{i-j}}{{\ell - j \choose i-j}} \leq \exp(r)$. Plugging in the polynomials $q_a(F)$ and rearranging yields
\[\ip{F,LF} \geq \delta(1 - \lambda_{r+1})(1- \gamma^{1/3}\exp(r)) - \sum_{j = 0}^r  \frac{c_j\ell^{j}}{\gamma}\E_{a \sim [n]^{j}}[q_a(F)(\delta(F^{ 2}|_a) - \gamma)] + B(F),\]
for some positive constants $c_j \leq \exp(r)$ and $B(F) = \frac{4}{3}\E_X[(F^{ 3} - F)\sum_{i = 0}^r F_i] + c\E_X[F - F^{ 4}]$, for a positive constant $c \leq \exp(r)\gamma^{1/3}$. 

We will now plug in upper bounds for $\lambda_i$. From Lemma~\ref{lem:eigenval}, we have that $\lambda_i = \frac{\binom{\ell - |T|}{(1 - \alpha)\ell - |T|}}{\binom{\ell}{(1 - \alpha)\ell}}$ for $i \leq (1-\alpha)\ell$ and $0$ otherwise. One can check that ${\lambda_i \leq (1-\alpha)\lambda_{i-1}}$ for all $i$ between $1$ and $\ell$. Since $\lambda_0 = 1$, we get that $\lambda_i \leq (1-\alpha)^i$, thus completing the proof of the lemma.

\end{proof}
